% ID: FIS-TER-EQ-002
% Titolo: Calorimetro con acqua e ghiaccio
% Disciplina: Fisica
% Area: Termologia
% Argomento: Temperatura di equilibrio
% Sottoargomento: Equilibrio termico con fusione del ghiaccio
% Classe: Seconda liceo scientifico
% Difficolta: 4
% Tipo: problema_numerico
% Risultato: soluzione da aggiungere
% Tempo_stimato: 12 min
% Competenze: bilancio termico, calore latente, scelta del modello fisico
% Prerequisiti: calore specifico, calore latente di fusione, equilibrio termico
% Tag: termologia, temperatura_equilibrio, calorimetria, ghiaccio, acqua, fusione
% Fonte: verifica_fisica_termologia_anonimizzata
% Autore: Riccardo Carli
% Licenza: CC-BY-SA-4.0

\begin{esercizio}
Un calorimetro contiene \(1{,}0\,\mathrm{kg}\) di acqua liquida a
\(28\,^\circ\mathrm{C}\) e \(100\,\mathrm{g}\) di ghiaccio a
\(-20\,^\circ\mathrm{C}\).

Trascurando la capacita termica del calorimetro e lo scambio di calore con
l'ambiente, determina la temperatura di equilibrio.
\end{esercizio}

\begin{soluzione}
% Soluzione da aggiungere.
\end{soluzione}

